\clearpage
\section{RX-side Noise Figure and Sensitivity}
\label{sec:rx-noise-figure}
\raggedbottom% this section's pages end naturally (no stretched gaps); restored at the end
% Generated by FR3_Characterization/plot_datasheet_nf.py -- do not edit by hand.
\newcommand{\nfLowMin}{5.2}
\newcommand{\nfLowMedian}{12.2}
\newcommand{\nfLowPninety}{16.5}
\newcommand{\nfLowMax}{22.9}
\newcommand{\nfLowCount}{224}
\newcommand{\nfHighMin}{5.6}
\newcommand{\nfHighMedian}{11.0}
\newcommand{\nfHighPninety}{15.8}
\newcommand{\nfHighMax}{20.5}
\newcommand{\nfHighCount}{203}
\newcommand{\nfMedian}{11.5}
\newcommand{\nfCount}{427}
\newcommand{\nfMinMargin}{14}
\newcommand{\nfOosCount}{15}
\newcommand{\nfOosMin}{10.4}
\newcommand{\nfOosMedian}{16.1}
\newcommand{\nfOosMax}{24.1}
\newcommand{\sensCW}{-162.5}
\newcommand{\sensOneM}{-102.5}
\newcommand{\sensFiveM}{-95.5}
\newcommand{\sensTenM}{-92.5}
\newcommand{\sensFiftyM}{-85.5}
\newcommand{\sensHundredM}{-82.5}
\newcommand{\sensFiveHundredM}{-75.5}
\newcommand{\nfStripeLow}{6.0}
\newcommand{\nfStripeHigh}{4.0}


This section first explains what noise figure is, how it sets the receiver's sensitivity, and how it was measured
(Section~\ref{sec:nf-concept}). It then gives the measured noise figure with the ``optimal'' AM3153 filter engaged
(Section~\ref{sec:nf-optimal}), the setting most users will run.

\subsection{Noise Figure: The Concept}
\label{sec:nf-concept}

\textbf{A Motivating Example.} Every receiver adds some noise of its own to whatever it receives. Even with
nothing at all connected to the ``RX-side FR3 RF In'' connector except a 50\,$\Omega$ termination, the termination's
own thermal noise comes in --- about −174\,dBm in every 1\,Hz of bandwidth at room temperature --- and every amplifier, mixer, and filter along the RX chain then adds more on top. The \emph{noise figure} (NF) measures how
much the receiver makes things worse: it is how many dB noisier the output is than it would be if the
receiver were a perfect, noiseless amplifier with the same gain. A receiver with NF~=~0\,dB adds nothing; one
with NF~=~6\,dB puts out four times the noise power a perfect receiver would. A lower NF means weaker signals
can still be heard above the noise.

\textbf{Which noise figure?} The mixer converts noise from \emph{both} RF frequencies that land on the same IF ---
the desired one and its image (see Section~\ref{sec:rx-image-rejection}). The NF in this section counts all of the
noise that reaches the IF output against the gain of the desired signal alone. That is the \emph{single-sideband}
(SSB) noise figure, the one that applies when receiving an ordinary signal that occupies only one side of the
LO.\footnote{Some datasheets quote a \emph{double-sideband} (DSB) NF instead, which credits the receiver for the
gain at both sidebands and therefore comes out up to 3\,dB lower. The SSB figure is the honest one for this unit's
normal use.}

\textbf{From noise figure to sensitivity.} \emph{Sensitivity} is the weakest signal the receiver can use,
and it follows directly from the NF:
\begin{center}
Sensitivity (dBm) = −174 + NF + 10·log\textsubscript{10}(bandwidth in Hz) + required SNR (dB)
\end{center}
The first three terms are the receiver's noise floor, referred to its input, over the signal's bandwidth;
the last is how far above that noise floor the signal has to be for the demodulator to work, which depends on
the modulation and coding, not on this unit. Every dB of NF is therefore a dB of sensitivity.

\textbf{How it's measured.} The unit's NF is found from two separate measurements at exactly the same
settings (same FR1, FR3, LO, and filter), both read at the ``RX-side FR1 IF Out'' connector:\footnote{This is the
``gain method''. It needs no calibrated noise source --- only a spectrum analyzer that is much quieter than
the unit's own output noise. Here the analyzer's own noise floor sat at least \nfMinMargin\,dB below the unit's
output noise at every point; it was measured separately and subtracted from every reading.}
\begin{enumerate}
  \item \textbf{Output noise.} With the RF input terminated in 50\,$\Omega$, measure the noise density at the IF
        output, in dBm/Hz.
  \item \textbf{Gain.} Feed in a weak CW tone of known power (−50\,dBm, well below the point where the chain
        starts to compress) and measure the IF output tone, to get the conversion gain.
\end{enumerate}
Both readings are referred to the unit's connectors: the losses of the cables between the unit and the analyzer
are removed from the output noise as well as from the gain. Then NF = (output noise) − (gain) + 174, all in dB. For example, an output noise of −120\,dBm/Hz with
45\,dB of gain gives NF = −120 − 45 + 174 = 9\,dB. The input tone power was itself measured, at every FR3
frequency, through the same cable used to feed the unit, rather than trusting the signal generator's
nominal setting.\footnote{The generator was within a few tenths of a dB of nominal across 6--26.5\,GHz,
with a few narrow dips of up to 0.7\,dB; the measured value is used throughout.} All readings use 100
power-domain averages.

\textbf{What's plotted.} NF was measured at \nfCount{} (FR1, FR3) combinations: FR1 from 1 to 4\,GHz and FR3
from 6 to 26.5\,GHz, both in 0.5\,GHz steps, for both low-side and high-side LO, wherever the LO falls between
7 and 22.5\,GHz, with the ``optimal'' filter engaged at every point. Grey cells
are combinations that would need the LO outside that range. In the heatmaps, \emph{darker is better} (lower NF),
the same ``darker is better'' sense as the gain figures, and both heatmaps share one colour scale so they can be
compared by eye. FR3 values above the unit's rated 24\,GHz maximum are measured and shown too, inside a red ``out of
design spec'' box. Because NF is a CW measurement, read at a single IF frequency, FR1 = 4\,GHz lies inside the
0.7--4.0\,GHz IF spec (unlike the conversion-gain figures, whose 500\,MHz windows run past it).

\textbf{A note about the 10\,MHz clock comb.} The noise is read 200\,kHz away from each IF frequency, not exactly on
it. The unit's 10\,MHz reference clock leaks a weak comb of spurious tones into the IF output, strongest at
multiples of 20\,MHz (with weaker ones at the odd multiples of 10\,MHz); it comes from the reference clock buffer,
and it sits a few dB to about 12\,dB above the noise floor at the strongest lines.
Reading 200\,kHz away keeps it out of the measurement. In use, it matters only when receiving an extremely weak CW
signal placed exactly on a multiple of 10\,MHz at the IF, where a comb line can sit on top of it; if that happens,
move the IF (or the LO) by a few hundred kHz. For a wideband signal it hardly matters: a 100\,MHz OFDM signal with
1024 subcarriers has only about ten subcarriers on comb lines, which together raise the total in-band noise by well
under 0.1\,dB.

\subsection{Noise Figure: Optimal Filter Engaged}
\label{sec:nf-optimal}

The AM3153 filter is what gives the RX side its image rejection (Section~\ref{sec:rx-image-rejection}), but it sits
after the LNA, so its insertion loss also adds to the noise figure. Here it is engaged at every point in its ``optimal''
setting, the same per-point choice used throughout this datasheet (Section~\ref{sec:sideband-suppression}).

\textbf{With the filter bypassed, the noise figure is lower.} Nothing then sits between the LNA and the mixer, so the
filter's loss no longer adds to the noise. The price is image rejection: a signal at the image frequency then reaches
the IF output almost unattenuated (Section~\ref{sec:rx-image-rejection}). Since most users engage the filter for its
image rejection, only the optimal-filter setting is characterized here. If your incoming FR3 signal is very weak and
nothing is transmitting at its image, bypassing the filter will lower the noise figure.

\textbf{Measurement conditions.} \lmxDriveRule{} (set from the LO frequency). Because the LO is shared, the TX side
was set up exactly as in normal use (its stored LO-leakage calibration applied and its optimal filter engaged), with
both TX-side ports terminated in 50\,\textohm. The unit was measured hot: no fan, LMX2820 die at
\condRxNfMinC--\condRxNfMaxC\textdegree C.

\textbf{Results.} Across all points up to the rated 24\,GHz, the median NF is \nfMedian\,dB. Over that same
range, with low-side LO the NF runs from \nfLowMin{} to \nfLowMax\,dB (median \nfLowMedian\,dB, and 90\% of
points at or below \nfLowPninety\,dB); with high-side LO it runs from \nfHighMin{} to \nfHighMax\,dB (median
\nfHighMedian\,dB, 90\% at or below \nfHighPninety\,dB) (Figures~\ref{fig:12-1}--\ref{fig:12-4}).
Above the rated 24\,GHz the unit still works, just out of design spec: the \nfOosCount{} points measured there (all
with low-side LO, boxed in red in Figure~\ref{fig:12-1}; the histograms count only the points within the design spec)
have an NF from \nfOosMin{} to \nfOosMax\,dB (median \nfOosMedian\,dB).

\textbf{Sensitivity.} At the median NF of \nfMedian\,dB, the input-referred noise floor (required SNR = 0\,dB) is:
\par\medskip\noindent\begin{minipage}{\linewidth}% caption and table stay on one page
\begin{center}
\captionof{table}{Input-referred noise floor at the median NF, optimal filter engaged.}\label{tab:nf-floor-optimal}
\vspace{0.4em}
\begin{tabular}{lc}
\toprule
Signal bandwidth & Noise floor at the RF input \\
\midrule
CW tone\textsuperscript{*} & \sensCW\,dBm/Hz \\
1\,MHz & \sensOneM\,dBm \\
5\,MHz & \sensFiveM\,dBm \\
10\,MHz & \sensTenM\,dBm \\
50\,MHz & \sensFiftyM\,dBm \\
100\,MHz & \sensHundredM\,dBm \\
500\,MHz & \sensFiveHundredM\,dBm \\
\bottomrule
\multicolumn{2}{l}{\footnotesize\textsuperscript{*}Per 1\,Hz of detection bandwidth (see below).}
\end{tabular}
\end{center}
\end{minipage}\par
Add the required SNR to get the sensitivity. For example, a signal needing 10\,dB of SNR in 100\,MHz has a
sensitivity 10\,dB above the 100\,MHz noise floor in the table. A CW tone has no bandwidth of its own, so its
noise floor depends on the bandwidth it is detected in: add 10·log\textsubscript{10} of that bandwidth in Hz to the CW row, e.g.\ +40\,dB for 10\,kHz.

\textbf{Where the NF is worst, and how to avoid it.} The clearest feature in both heatmaps is a stripe of
higher NF wherever the LO sits at 14.0 or 14.5\,GHz: there the median NF is about \nfStripeLow\,dB (low-side LO)
and \nfStripeHigh\,dB (high-side LO) worse than with the LO at the neighbouring 13--13.5 and 15--15.5\,GHz, at every
FR1. Because it follows the LO frequency rather than the RF frequency, it shows up as a diagonal stripe in the
low-side heatmap (FR3 = LO + FR1) and as the mirror diagonal in the high-side one (FR3 = LO − FR1). If your FR3 and
FR1 put the LO at 14 to 14.5\,GHz, switching to the other LO side moves the LO by twice the FR1 frequency, well
clear of that stripe, usually for a much better NF. None of this matters unless the incoming FR3 signal is weak
enough to sit close to the receiver's noise floor: a stronger signal stays far above the added noise either way.

\medskip
\figpair{images/sec12_fig1.pdf}{Noise figure, low-side LO: full measured (FR1, FR3) grid, optimal filter.
Zoom in to read each cell's exact values, printed as (FR1, FR3, noise figure), with FR1 and FR3 in GHz and the
noise figure in dB.}{fig:12-1}%
        {images/sec12_fig2.pdf}{LO: Low-side. Filter: Optimal. Distribution of Figure~\ref{fig:12-1}, points within the design spec.}{fig:12-2}
\figpair{images/sec12_fig3.pdf}{Noise figure, high-side LO: full measured (FR1, FR3) grid, optimal filter.
Zoom in to read each cell's exact values, printed as (FR1, FR3, noise figure), with FR1 and FR3 in GHz and the
noise figure in dB.}{fig:12-3}%
        {images/sec12_fig4.pdf}{LO: High-side. Filter: Optimal. Distribution of Figure~\ref{fig:12-3}, points within the design spec.}{fig:12-4}
\clearpage
\flushbottom
